\documentclass[12pt]{article}
\usepackage[left=0.75in,right=0.75in,top=0.75in,bottom=0.55in]{geometry}
\usepackage{fontspec}
\setmainfont{Times New Roman}
\usepackage[hidelinks]{hyperref}
\usepackage{parskip}
\usepackage{setspace}
\setlength{\parskip}{0.35em}
\pagestyle{empty}

\begin{document}
\onehalfspacing

Dear Scale AI Hiring Team,

I am applying for the Security Engineer, Product Security position at Scale AI. I bring more than ten years of software engineering and cybersecurity experience across critical-infrastructure and public-sector environments. My work combines software development with threat modeling, vulnerability assessment, penetration testing, red teaming, secure coding remediation, malware and APT analysis, incident response, and technical reporting. Scale's work developing dependable AI systems for consequential decisions closely matches my commitment to responsible security engineering.

I began my career at KEPCO E\&C developing reactor design software and engineering databases for nuclear-plant workflows while performing vulnerability assessments and secure coding remediation for critical information and communications infrastructure. This work showed me how product design, implementation choices, operational constraints, and security controls interact in high-consequence environments. It also established the approach I continue to use: understand how a system is designed and used, identify realistic attack paths, explain their impact clearly, and develop mitigations that engineering teams can sustain.

At the Ministry of Science and ICT, I expanded that foundation across public-sector and critical-infrastructure environments. I delivered penetration testing, red teaming, vulnerability analysis, and remediation consulting to more than 200 public institutions, including nuclear, thermal-power, railway-control, and power-system environments. In my current threat-management work, I conduct malware analysis, attack-evidence collection, APT tracking, technical reporting, and remediation guidance, contributing more than 500 public cyber-threat intelligence cases each year. These responsibilities demand independent investigation, root-cause analysis, and clear communication with engineers and decision-makers.

My interest in securing AI-enabled systems is grounded in practical security work. In 2023, I initiated a research project with the Korea Internet \& Security Agency on cyber incident response for critical infrastructure. We analyzed hacking infrastructure used by international threat groups and explored how AI and data analytics could improve attack prediction, detection, and mitigation. At the University of California, Davis, I am building on that experience through graduate study of machine learning, large language models, generative AI, RAG, model evaluation, security automation, and MLOps.

Scale's need for product-security engineering across the AI/ML software ecosystem strongly matches this background. I can contribute experience reviewing software and system risk, conducting hands-on assessments, guiding secure coding remediation, and communicating vulnerability mechanics, exploitability, and impact. My programming background in Python, Java, C/C++, C\#, JavaScript, and SQL helps me work across security analysis and engineering implementation. I would welcome the opportunity to help Scale strengthen security throughout product design and development while deepening my expertise in cloud-native security automation, CI/CD controls, and infrastructure orchestration.

Thank you for your consideration.

Sincerely,\\
Kevin Lee

\end{document}
